

\paragraph{Terminology.}
\subparagraph{(A) Shape dimension (\(\ShapeRank\)).}
The number of \(r\)-dependent patterns (columns of \(\Phi\)) needed, after removing rigid modes, to span the Saint–Venant residual over the section. 
Formally, it is the \emph{separation rank} of the residual
\(R(x,\bm{r})\in \mathcal F_x\otimes(\mathcal V_r/\mathscr R)\): the smallest \(N\) such that
\[
R(x,\bm{r})=\sum_{j}^N \alpha_j(\reflenx)\,\bm{\Phi}_j(\bm{r})
\quad\text{with}\quad \bm{\Phi}_j\ \text{independent of }\reflenx.
\]
For a generic section: \(\ShapeRank=8\) (or \(9\) if the uniform axial \(\reflenx^2\) term is kept as a ``shape").

From the matrix–vector viewpoint \(\bm\Phi(\bm{r})\bm\alpha(\reflenx)\), take \(\bm\Phi(\bm{r})\) with \(\ShapeRank\) columns (the fixed warping shapes \(\bm{\Phi}_j\) after quotienting rigid modes).
The \emph{shape rank} is the minimal such number of columns.

\subparagraph{(B) Amplitude dimension (\(\FieldRank\); free‑amplitude rank).}
The number of independent scalar functions of \(\reflenx\) required \emph{after} associating as many amplitude channels as possible to the 1D measures \((\epsilon,\bm\gamma,\bm\kappa)\).
Concretely, if a shape’s required amplitude \(\alpha_j(\reflenx)\) has exactly the same \(\reflenx\)-law as a component (or a fixed linear combination) of \(\epsilon,\gamma,\kappa\), set \(\alpha_j(\reflenx)\) equal to that measure—then it does not count as a new unknown function.
What remains are the \emph{unavoidable} amplitude functions. For a generic section: \(\FieldRank=5\).

From the matrix–vector viewpoint, \(\bm\alpha(\reflenx)\in\mathbb R^{\ShapeRank}\) collects the channels \(\alpha_j(\reflenx)\).
Impose identifications of the form
\[
\alpha_j(\reflenx)=\mathcal L_j\big[\epsilon(\reflenx),\bm\gamma(\reflenx),\bm\kappa(\reflenx)\big]
\]
with \(\mathcal L_j\) a fixed, geometry‑dependent linear mapping that works for \emph{all} load parameters.
Every channel admitting such an identification is not an independent unknown.
The \emph{free‑amplitude rank} \(\FieldRank\) is the number of entries of \(\bm\alpha(\reflenx)\) that cannot be written this way.

\medskip
These numbers differ because many shapes can share the same driving signal in \(\reflenx\).
Shapes control \emph{where} the field lives across the section; amplitudes control \emph{how} it varies along \(\reflenx\).
Hence one often needs more distinct shapes than distinct \(\reflenx\)-signals.


\[
\begin{aligned}
\ShapeRank&=\operatorname{rank}_{\mathrm{sep}} R
\quad\text{(minimal number of \(\reflenx\)-independent non‑rigid shapes)}\\
\FieldRank&=\dim \mathcal G/\mathcal E=\ShapeRank-r_{\mathrm{assoc}}
\end{aligned}
\]
Here \(\mathcal G=\mathrm{span}\{\alpha_j(\reflenx)\}\) and \(\mathcal E\) is the subspace generated by fixed linear combinations of \((\epsilon,\gamma,\kappa)\);
for a generic section \(r_{\mathrm{assoc}}=3\) (Poisson extension \(1\) + flexural axial pair \(2\)), so \(\FieldRank=8-3=5\).

\paragraph{Rigid Quotient}

We obtain clean bounds by treating the SV field as a separated function of \(\reflenx\) and \(\bm{r}\), then quotienting out what the kinematics \((\bm x,\bm\theta)\) can generate.

\begin{proposition}\label{prop:min8}
For a generic cross–section (no symmetry assumptions, arbitrary \(\CentroidLocation\)), the full Saint–Venant family
\((\bm a_{\Section},a_\varepsilon,a_\vartheta,\bm b_{\Section})\) admits
\[
\boxed{\ \text{exactly } \FieldRank = 5\ \text{independent warping amplitude functions of }\reflenx\ }
\]
after associating as many amplitudes as possible to \(\bm\gamma(\reflenx),\bm\kappa(\reflenx)\).
The mathematically minimal number of \(\reflenx\)-independent warping shapes is
\[
\boxed{\ \ShapeRank=8\ }
\]
(or \(9\) if you refuse to absorb the uniform axial \(\reflenx^2\) term into \(\bm u\)).
\end{proposition}

\subparagraph{Rigid span and residual.}
Let
\[
\mathscr{R}(\Section)
=\operatorname{span}\Big(\underbrace{\{\PlaneBasis_y,\PlaneBasis_z,\FrameBasis\times\bm{r}\}}_{\text{transverse}}\ \oplus\ 
\underbrace{\{\FrameBasis\otimes\bm{r},\FrameBasis\}}_{\text{axial}}\Big).
\]
Project the SV field onto \(\mathscr R(\Section)\) (choose \(\bm u,\bm\theta\) to cancel those parts) and call the residual
\[
\mathbf{R}(x,\bm{r}) \in \mathcal{F}_x \otimes \big(\mathcal{V}_r / \mathscr{R}(\Section)\big),
\qquad \mathcal{F}_x=\mathrm{span}\{1,x,\reflenx^2,\reflenx^3\}.
\]
The minimum number of \emph{warping shapes} required is the separation rank of \(\mathbf R\):
\[
\ShapeRank=\operatorname{rank}_\text{sep}\mathbf R,
\]
i.e., the smallest \(N\) so that \(\mathbf R(x,\bm{r})=\sum_{j}^N \alpha_j(\reflenx)\,\bm W_j(\bm{r})\) with each \(\
\bm W_j\) independent of \(\reflenx\).

\subparagraph{Eight non‑rigid shapes (generic).}
After removing rigids shapes, the SV field leaves the \(\reflenx\)-independent shapes
\[
\big\{\,W_{(\varepsilon)};\ W^{(a)}_y,W^{(a)}_z;\ W^{(b)}_y,W^{(b)}_z;\ \FrameBasis\!\otimes\!\psi_y,\ \FrameBasis\!\otimes\!\psi_z;\ \FrameBasis\,\varphi\,\big\},
\]
giving \(\ShapeRank=8\) (or \(9\) if you keep the uniform axial \(\reflenx^2\) term as a shape).

\subparagraph{Proof sketch for the counts.}
Enumerate the \(\reflenx\)-monomials, assemble the corresponding \(r\)-shapes, project out \(\mathscr R(\Section)\), and take the span in \(\mathcal{V}_r/\mathscr R\) to obtain \(\ShapeRank\).
Then associate Poisson extension and the flexural axial pair to \((\epsilon,\bm\gamma)\), leaving exactly five unavoidable amplitude functions (two \(a\)-Poisson, two \(b\)-Poisson, and one torsion—or, in another gauge, the uniform axial \(\reflenx^2\) term).

\subparagraph{General principle.}
\[
\boxed{\ \FieldRank \le \ShapeRank,\ \ \text{with strict inequality whenever multiple shapes share the same \(\reflenx\)-driver from }(\epsilon,\gamma,\kappa). \ }
\]
